\documentclass[10pt,a4paper]{amsart}
\usepackage{amsmath,amssymb,amsthm,graphicx}
\usepackage{verbatim}
\newenvironment{reference}{\comment}{\endcomment}
%\newenvironment{reference}{}{}
\newenvironment{slogan}{\comment}{\endcomment}
\newenvironment{history}{\comment}{\endcomment}

% For commutative diagrams we use Xy-pic
\usepackage[all]{xy}

% We use 2cell for 2-commutative diagrams.
\xyoption{2cell}
% Theorem environments.
%
\theoremstyle{plain}
\newtheorem{theorem}[subsection]{Theorem}
\newtheorem{proposition}[subsection]{Proposition}
\newtheorem{lemma}[subsection]{Lemma}

\theoremstyle{definition}
\newtheorem{definition}[subsection]{Definition}
\newtheorem{example}[subsection]{Example}
\newtheorem{exercise}[subsection]{Exercise}
\newtheorem{situation}[subsection]{Situation}

\theoremstyle{remark}
\newtheorem{remark}[subsection]{Remark}
\newtheorem{remarks}[subsection]{Remarks}

\numberwithin{equation}{subsection}

% Macros
%
\def\lim{\mathop{\mathrm{lim}}\nolimits}
\def\colim{\mathop{\mathrm{colim}}\nolimits}
\def\Spec{\mathop{\mathrm{Spec}}}
\def\Hom{\mathop{\mathrm{Hom}}\nolimits}
\def\Ext{\mathop{\mathrm{Ext}}\nolimits}
\def\SheafHom{\mathop{\mathcal{H}\!\mathit{om}}\nolimits}
\def\SheafExt{\mathop{\mathcal{E}\!\mathit{xt}}\nolimits}
\def\Sch{\mathit{Sch}}
\def\Mor{\mathop{\mathrm{Mor}}\nolimits}
\def\Ob{\mathop{\mathrm{Ob}}\nolimits}
\def\Sh{\mathop{\mathit{Sh}}\nolimits}
\def\NL{\mathop{N\!L}\nolimits}
\def\CH{\mathop{\mathrm{CH}}\nolimits}
\def\proetale{{pro\text{-}\acute{e}tale}}
\def\etale{{\acute{e}tale}}
\def\QCoh{\mathit{QCoh}}
\def\Ker{\mathop{\mathrm{Ker}}}
\def\Im{\mathop{\mathrm{Im}}}
\def\Coker{\mathop{\mathrm{Coker}}}
\def\Coim{\mathop{\mathrm{Coim}}}

% Boxtimes
%
\DeclareMathSymbol{\boxtimes}{\mathbin}{AMSa}{"02}

%
% Macros for moduli stacks/spaces
%
\def\QCohstack{\mathcal{QC}\!\mathit{oh}}
\def\Cohstack{\mathcal{C}\!\mathit{oh}}
\def\Spacesstack{\mathcal{S}\!\mathit{paces}}
\def\Quotfunctor{\mathrm{Quot}}
\def\Hilbfunctor{\mathrm{Hilb}}
\def\Curvesstack{\mathcal{C}\!\mathit{urves}}
\def\Polarizedstack{\mathcal{P}\!\mathit{olarized}}
\def\Complexesstack{\mathcal{C}\!\mathit{omplexes}}
% \Pic is the operator that assigns to X its picard group, usage \Pic(X)
% \Picardstack_{X/B} denotes the Picard stack of X over B
% \Picardfunctor_{X/B} denotes the Picard functor of X over B
\def\Pic{\mathop{\mathrm{Pic}}\nolimits}
\def\Picardstack{\mathcal{P}\!\mathit{ic}}
\def\Picardfunctor{\mathrm{Pic}}
\def\Deformationcategory{\mathcal{D}\!\mathit{ef}}
\usepackage{cleveref}
\crefname{originalreference}{original reference}{original references}
\Crefname{originalreference}{Original reference}{Original references}
\makeatletter
\newlabel{sites-example-indiscrete}{{07GE}{source}{}{sourcepointer-1}{}}
\newlabel{sites-example-indiscrete@cref}{{[originalreference][0][]07GE}{[1][0][]source}{}{}{}}
\makeatother
\crefname{originalreference}{original reference}{original references}
\Crefname{originalreference}{Original reference}{Original references}
\makeatletter
\newlabel{homology-definition-serre-subcategory}{{02MO}{source}{}{sourcepointer-2}{}}
\newlabel{homology-definition-serre-subcategory@cref}{{[originalreference][0][]02MO}{[1][0][]source}{}{}{}}
\makeatother
\crefname{originallemma}{original lemma}{original lemmas}
\Crefname{originallemma}{Original lemma}{Original lemmas}
\makeatletter
\newlabel{lemma-chaotic-quasi-coherent}{{0GZN}{source}{}{sourcepointer-3}{}}
\newlabel{lemma-chaotic-quasi-coherent@cref}{{[originallemma][0][]0GZN}{[1][0][]source}{}{}{}}
\makeatother
\crefname{originalreference}{original reference}{original references}
\Crefname{originalreference}{Original reference}{Original references}
\makeatletter
\newlabel{homology-lemma-five-lemma}{{05QB}{source}{}{sourcepointer-4}{}}
\newlabel{homology-lemma-five-lemma@cref}{{[originalreference][0][]05QB}{[1][0][]source}{}{}{}}
\makeatother
\title{Quasi-coherent modules on a chaotic site with flat restrictions form a weak Serre subcategory}
\author{The Stacks Project Authors}
\date{stacks-project:sites-modules.tex:local-snapshot-20261012; GNU FDL 1.2 or later}
\begin{document}
\maketitle
\noindent\textit{Editorial note.} All delivered mathematical text is the author's own, copied verbatim from the pinned Stacks Project chapter file sites-modules.tex; no statement or proof step was rewritten, reordered or removed.\par\medskip
\noindent\textit{Editorial note.} The proof contains a commutative diagram typeset with Xy-pic; the retained two-slice preamble of this source supplies the native Xy-pic setup, so it compiles as in the source.\par\medskip
\noindent\textit{Editorial note.} The counted span cites two external items (an example and a definition), recorded in assembly.external\_references with canonical tags as navigation pointers.\par\medskip
\noindent\textit{Allowed standard prerequisites (editorial).}\begin{itemize}
\item Characterization of quasi-coherence in the chaotic topology by base change. Stacks Project, Modules on Sites, Lemma lemma-chaotic-quasi-coherent (tag 09BA). Turns quasi-coherence into the module-theoretic base-change condition that the diagram chase tests.
\item Weak Serre subcategories and the five lemma. Stacks Project, Homology, Definition homology-definition-serre-subcategory (tag 02MP). Supplies the property to be verified and the tool for the middle-term conclusion.
\end{itemize}\medskip
\ifdefined\hypertarget\hypertarget{sourcepointer-1}{}\fi
\noindent\textit{Source reference.} Original reference 07GE: sources/stacks-chow/tags.txt, source lines 9514--9514. Cross-chapter standard prerequisite cited by the author inside the counted span; it is used as a pointer and its own proof is not part of this item.\par\medskip
\ifdefined\hypertarget\hypertarget{sourcepointer-2}{}\fi
\noindent\textit{Source reference.} Original reference 02MO: sources/stacks-chow/tags.txt, source lines 3451--3451. Cross-chapter standard prerequisite cited by the author inside the counted span; it is used as a pointer and its own proof is not part of this item.\par\medskip
\ifdefined\hypertarget\hypertarget{sourcepointer-3}{}\fi
\noindent\textit{Source reference.} Original lemma 0GZN: sources/stacks-chow/tags.txt, source lines 21183--21183. Cross-chapter standard prerequisite cited by the author inside the counted span; it is used as a pointer and its own proof is not part of this item.\par\medskip
\ifdefined\hypertarget\hypertarget{sourcepointer-4}{}\fi
\noindent\textit{Source reference.} Original reference 05QB: sources/stacks-chow/tags.txt, source lines 7370--7370. Cross-chapter standard prerequisite cited by the author inside the counted span; it is used as a pointer and its own proof is not part of this item.\par\medskip
\begin{lemma}
\label{lemma-chaotic-flat-quasi-coherent}
Let $\mathcal{C}$ be a category viewed as a site with the chaotic
topology, see Sites, Example \ref{sites-example-indiscrete}.
Let $\mathcal{O}$ be a sheaf of rings on $\mathcal{C}$.
Assume for all $U \to V$ in $\mathcal{C}$ the restriction map
$\mathcal{O}(V) \to \mathcal{O}(U)$ is a flat ring map.
Then the category of quasi-coherent $\mathcal{O}$-modules is
a weak Serre subcategory of $\textit{Mod}(\mathcal{O})$.
\end{lemma}

\begin{proof}
We will check the definition of a weak Serre subcategory, see
Homology, Definition \ref{homology-definition-serre-subcategory}.
To do this we will use the characterization of quasi-coherent modules given in
Lemma \ref{lemma-chaotic-quasi-coherent}.
Consider an exact sequence
$$
\mathcal{F}_0 \to \mathcal{F}_1 \to
\mathcal{F}_2 \to \mathcal{F}_3 \to
\mathcal{F}_4
$$
in $\textit{Mod}(\mathcal{O})$ with $\mathcal{F}_0$, $\mathcal{F}_1$,
$\mathcal{F}_3$, and $\mathcal{F}_4$ quasi-coherent. Let $U \to V$
be a morphism of $\mathcal{C}$ and consider the commutative diagram
$$
\xymatrix{
\mathcal{F}_0(V) \otimes_{\mathcal{O}(V)} \mathcal{O}(U) \ar[r] \ar[d] &
\mathcal{F}_1(V) \otimes_{\mathcal{O}(V)} \mathcal{O}(U) \ar[r] \ar[d] &
\mathcal{F}_2(V) \otimes_{\mathcal{O}(V)} \mathcal{O}(U) \ar[r] \ar[d] &
\mathcal{F}_3(V) \otimes_{\mathcal{O}(V)} \mathcal{O}(U) \ar[r] \ar[d] &
\mathcal{F}_4(V) \otimes_{\mathcal{O}(V)} \mathcal{O}(U) \ar[d] \\
\mathcal{F}_0(U) \ar[r] &
\mathcal{F}_1(U) \ar[r] &
\mathcal{F}_2(U) \ar[r] &
\mathcal{F}_3(U) \ar[r] &
\mathcal{F}_4(U)
}
$$
By assumption the vertical arrows with indices $0$, $1$, $3$, $4$
are isomorphisms. Since the topology on $\mathcal{C}$ is chaotic
taking sections over an object of $\mathcal{C}$ is exact and hence
the lower row is exact. Since $\mathcal{O}(V) \to \mathcal{O}(U)$
is flat also the upper row is exact. Thus we conclude that the middle
arrow is an isomorphism by the $5$ lemma
(Homology, Lemma \ref{homology-lemma-five-lemma}).
\end{proof}

\end{document}
